\begin{frame}{DSU}
    \metroset{block=fill}

    No DSU, podemos fazer uma segunda otimização. Ao procurar um representante de um elemento, o faremos apontar diretamente para a raiz, para que próximas consultas sejam melhores. O algoritmo para encontrar o representante de um elemento $u$ é:

    \begin{itemize}
        \item Se $u = pai[u]$, então retornamos u.
        \item Se $u \not = $pai$[u]$, então recursivamente encontramos o representante do pai de $u$, e definimos o pai de $u$ como este representante.
    \end{itemize}
    
    Pode ser provado que a complexidade final, com esta otimização, é $O((q + n)\alpha(n))$, onde $\alpha(n)$ é a inversa da função de Ackermann. Na prática, podemos considerar que a complexidade é $O(q+n)$.

    % \begin{block}{Implementação da busca binária}
    %     \footnotesize
    %     \inputminted{cpp}{codigos/dsu.cpp}    
    % \end{block}


\end{frame}
